\documentclass[a4paper]{article}

% Paquetes básicos
% Español (México) con XeLaTeX: fontspec para fuentes Unicode, polyglossia
% para la división silábica y los nombres del documento en la variante mexicana.
\usepackage{fontspec}
\usepackage{polyglossia}
\setdefaultlanguage[variant=mexican]{spanish}
% Los nombres chinos van como «pinyin (original)»: xeCJK da a los caracteres
% Han su propia fuente; sin ella XeLaTeX los omite con un aviso.
% FandolSong no cubre algunos caracteres de nombres (U+73FA); WenQuanYi Zen Hei sí.
\usepackage[AutoFallBack=true]{xeCJK}
\setCJKmainfont{FandolSong-Regular.otf}
\setCJKfallbackfamilyfont{\CJKrmdefault}{WenQuanYi Zen Hei}
% polyglossia no traduce los nombres de listados.
\AtBeginDocument{\providecommand{\lstlistingname}{}\renewcommand{\lstlistingname}{Listado}%
  \providecommand{\lstlistlistingname}{}\renewcommand{\lstlistlistingname}{Índice de listados}}
\usepackage{amsmath, amssymb}
\usepackage{graphicx}
\usepackage[margin=2.5cm]{geometry}
\usepackage[most]{tcolorbox}
\usepackage{etoolbox}
\usepackage{listings}
\usepackage{hyperref}
\usepackage{booktabs}
\usepackage{subcaption}
\usepackage{float}

% Cajas de resaltado
\newtcolorbox{knowledgebox}[1]{
    enhanced, colback=blue!5!white, colframe=blue!75!black, colbacktitle=blue!75!black,
    coltitle=white, fonttitle=\bfseries, title={#1}, attach boxed title to top left={yshift=-2mm, xshift=2mm},
    boxrule=1pt, sharp corners
}

\newtcolorbox{importantbox}[1]{
    enhanced, colback=yellow!10!white, colframe=yellow!80!black, colbacktitle=yellow!80!black,
    coltitle=black, fonttitle=\bfseries, title={#1}, sharp corners
}

\newtcolorbox{warningbox}[1]{
    enhanced, colback=red!5!white, colframe=red!75!black, colbacktitle=red!75!black,
    coltitle=white, fonttitle=\bfseries, title={#1}, sharp corners
}

% Listados
\lstset{
    language=Python,
    basicstyle=\ttfamily\small,
    keywordstyle=\color{blue},
    stringstyle=\color{red!60!black},
    commentstyle=\color{green!60!black},
    breaklines=true,
    frame=single,
    numbers=left,
    numberstyle=\tiny\color{gray},
    captionpos=b,
    extendedchars=true
}

% Metadatos de la página inicial
\newcommand{\notetitle}{KAIST CS492D: Modelos de difusión en tiempo continuo}
\newcommand{\noteauthors}{Elaboradas a partir de la clase de Minhyuk Sung}
\newcommand{\notedate}{Otoño de 2025}
\newcommand{\videochannel}{Minhyuk Sung (KAIST)}
\newcommand{\videopublishdate}{2025}
\newcommand{\videoduration}{55 minutos}
\newcommand{\videourl}{https://www.youtube.com/watch?v=by_flhcQJj4}
\newcommand{\videocoverpath}{cover.jpg}
\newcommand{\repourl}{https://github.com/hqhq1025/ai-course-notes}

\begin{document}

\begin{titlepage}
\centering
{\Large Notas del curso\par}
\vspace{1.2cm}
{\huge\bfseries \notetitle\par}
\vspace{0.8cm}
{\large \noteauthors\par}
\vspace{0.3cm}
{\large \notedate\par}
\vspace{1.2cm}

\ifdefempty{\videocoverpath}{
{\small Al generar las notas, indique la ruta local de la portada del video.\par}
}{
\includegraphics[width=0.82\textwidth,height=0.45\textheight,keepaspectratio]{\videocoverpath}\par
}

\vfill
\begin{tcolorbox}[width=0.9\textwidth, colback=black!2!white, colframe=black!60, sharp corners]
\textbf{Autor o canal del video}: \videochannel\par
\textbf{Fecha de publicación}: \videopublishdate\par
\textbf{Duración del video}: \videoduration\par
\ifdefempty{\videourl}{}{
\textbf{Enlace del video}: \href{\videourl}{\nolinkurl{\videourl}}\par
}
\end{tcolorbox}
\end{titlepage}

\tableofcontents
\newpage

%% --- Inicio del contenido --- %%

\section{Introducción: de lo discreto a lo continuo}

En las clases anteriores, hemos estudiado los fundamentos de los modelos de difusión (Diffusion Models), incluyendo el proceso hacia adelante y el proceso hacia atrás del DDPM (Denoising Diffusion Probabilistic Models), así como el método del DDIM (Denoising Diffusion Implicit Models) para convertir un proceso estocástico en uno determinista con el fin de reducir el número de pasos de muestreo. Estos conocimientos básicos son suficientes para construir y entrenar sus propios modelos de difusión en muchas aplicaciones prácticas.

El objetivo central de esta clase es comprender los modelos de difusión desde una nueva perspectiva: el \textbf{dominio continuo del tiempo}. Veremos que:

\begin{itemize}
    \item Tanto el DDPM de tiempo discreto como el SMLD (Score Matching con Dinámica de Langevin) pueden unificarse bajo un marco de Ecuaciones Diferenciales Estocásticas (SDE, Stochastic Differential Equation);
    \item Dada una SDE hacia adelante, se puede derivar directamente la SDE hacia atrás;
    \item Existe una forma equivalente a la SDE en términos de Ecuaciones Diferenciales Ordinarias (ODE, Ordinary Differential Equation), denominada ODE de flujo de probabilidad;
    \item Este marco continuo nos otorga mayor libertad y herramientas teóricas para diseñar modelos de difusión.
\end{itemize}

\begin{knowledgebox}{Antecedentes del curso}
Esta clase se basa en el artículo clásico de Song et al. (2021) titulado ``Score-Based Generative Modeling through Stochastic Differential Equations''. Este trabajo unifica los dos paradigmas SMLD (NCSN) y DDPM, proponiendo un marco general de SDE. Es una literatura fundamental para comprender la teoría moderna de los modelos de difusión.
\end{knowledgebox}

\begin{importantbox}{Puntos clave}
Anteriormente definimos el proceso hacia adelante y hacia atrás de los modelos de difusión bajo \textbf{tiempo discreto}. El cambio clave en esta clase es que, cuando el número de pasos discretos $N \to \infty$, todo el proceso de adición de ruido hacia adelante puede describirse mediante una SDE, mientras que el proceso de eliminación de ruido hacia atrás corresponde igualmente a una SDE inversa. Esta perspectiva de tiempo continuo no solo unifica el SMLD y el DDPM, sino que también proporciona la base teórica para diseñar nuevos modelos de difusión y muestreadores eficientes.
\end{importantbox}
\section{Fundamentos de las ecuaciones diferenciales estocásticas (SDE)}

\subsection{Forma general de las SDE}

En el dominio de tiempo continuo, el proceso de difusión hacia adelante de los modelos de difusión puede representarse mediante la siguiente SDE:

$$
\mathrm{d}\mathbf{x} = f(\mathbf{x}, t)\,\mathrm{d}t + g(t)\,\mathrm{d}\mathbf{w}
$$

Donde cada símbolo tiene el siguiente significado:

\begin{itemize}
    \item $\mathbf{x}$ —— Variable de datos (por ejemplo, imágenes), que es un proceso estocástico en función del tiempo $t$;
    \item $f(\mathbf{x}, t)$ —— \textbf{coeficiente de deriva} (drift coefficient), que determina la dirección de evolución temporal de la media de $\mathbf{x}_t$;
    \item $g(t)$ —— \textbf{coeficiente de difusión} (diffusion coefficient), que controla la intensidad del ruido inyectado;
    \item $\mathbf{w}$ —— \textbf{proceso de Wiener estándar} (proceso de Wiener), también conocido como movimiento browniano;
    \item $\mathrm{d}\mathbf{w}$ —- incremento infinitesimal del proceso de Wiener.
\end{itemize}

\begin{importantbox}{Intuición sobre la deriva y la difusión}
El \textbf{término de deriva} $f(\mathbf{x}, t)\,\mathrm{d}t$ puede entenderse como una "fuerza" determinista que mueve cada punto de datos en cierta dirección (por ejemplo, contrayéndose hacia el origen). El \textbf{término de difusión} $g(t)\,\mathrm{d}\mathbf{w}$ inyecta ruido aleatorio, haciendo que la trayectoria de $\mathbf{x}_t$ se vuelva aserrada. La combinación de ambos produce un mapeo continuo desde la distribución de datos hasta la distribución de ruido.
\end{importantbox}

\subsection{Proceso de Wiener (movimiento browniano)}

El proceso de Wiener $\mathbf{w}(t)$ es la fuente de aleatoriedad en las SDE y posee las siguientes cuatro propiedades características:

\begin{enumerate}
    \item \textbf{Condición inicial}: $\mathbf{w}(0) = 0$ (casi seguro, es decir, $P(\mathbf{w}(0)=0) = 1$);
    \item \textbf{Normalidad de los incrementos}: para cualquier $0 \le s < t$, el incremento $\mathbf{w}(t) - \mathbf{w}(s) \sim \mathcal{N}(0, (t-s)\mathbf{I})$;
    \item \textbf{Incrementos independientes}: los incrementos en intervalos de tiempo no superpuestos son mutuamente independientes;
    \item \textbf{Continuidad de las trayectorias}: $\mathbf{w}(t)$ es continuo respecto a $t$ casi seguro.
\end{enumerate}

\begin{knowledgebox}{El propio proceso de Wiener es una SDE}
Si se establece $f = 0$ y $g = 1$, la SDE se reduce a $\mathrm{d}\mathbf{x} = \mathrm{d}\mathbf{w}$, cuya solución es el propio proceso de Wiener. Por lo tanto, el movimiento browniano puede considerarse la SDE más simple: no hay deriva, solo difusión puramente aleatoria.
\end{knowledgebox}

\subsection{Desarrollo en serie de Taylor y aproximación}

Al convertir el proceso de difusión discreto a una SDE continua, el desarrollo en serie de Taylor es la herramienta matemática fundamental. Dada una función suave $f$, su aproximación de primer orden es:

$$
f(t + \Delta t) \approx f(t) + f'(t) \cdot \Delta t
$$

Por ejemplo, para $f(\epsilon) = \sqrt{1 - \epsilon}$, cuando $\epsilon$ es muy pequeño:

$$
\sqrt{1 - \epsilon} \approx 1 - \frac{\epsilon}{2}
$$

Esta aproximación se utiliza repetidamente al deducir las funciones $f(t)$ y $g(t)$ del VP-SDE. Cuando el tamaño del paso temporal $\Delta t \to 0$, los términos de orden superior desaparecen, obteniendo así la forma exacta de la ecuación diferencial.

\begin{warningbox}{Supuestos previos a la aproximación continua}
La conversión de lo discreto a lo continuo depende de las hipótesis de que el tamaño del paso temporal es suficientemente pequeño ($\Delta t \to 0$) y que la función de programación del ruido (como $\beta_t$) es suficientemente suave. Para la discretización con un número finito de pasos en implementaciones prácticas, la SDE continua es solo una aproximación: a mayor número de pasos, la aproximación será más precisa. Esto también explica por qué aumentar el número de pasos de discretización suele mejorar la calidad generada.
\end{warningbox}

\subsection{Resumen de la sección}

Las SDE describen la evolución aleatoria de los datos en tiempo continuo a través de los términos de deriva y difusión. El proceso de Wiener proporciona el impulso estocástico estándar. Mediante el desarrollo en serie de Taylor, podemos partir de las probabilidades de transición de una cadena de Markov discreta y deducir la SDE continua correspondiente.
\section{SDE de varianza explosiva: VE-SDE}

\subsection{De SMLD a VE-SDE}

SMLD (Score Matching con Dinámica de Langevin, también conocido como NCSN) define la distribución de transición hacia adelante como:

$$
q(\mathbf{x}_t \mid \mathbf{x}_{t-1}) = \mathcal{N}\!\left(\mathbf{x}_t;\, \mathbf{x}_{t-1},\, (\sigma_t^2 - \sigma_{t-1}^2)\mathbf{I}\right)
$$

La distribución de salto correspondiente (de $\mathbf{x}_0$ directamente a $\mathbf{x}_t$) es:

$$
q(\mathbf{x}_t \mid \mathbf{x}_0) = \mathcal{N}\!\left(\mathbf{x}_t;\, \mathbf{x}_0,\, \sigma_t^2 \mathbf{I}\right)
$$

\begin{importantbox}{La idea central de SMLD}
El proceso hacia adelante en SMLD consiste esencialmente en \textbf{incrementar gradualmente la varianza del ruido gaussiano}, sin desplazar la media de los puntos de datos. Es decir, cada punto de datos $\mathbf{x}_0$ se "difumina": a medida que avanza el tiempo, la distribución gaussiana centrada en $\mathbf{x}_0$ se vuelve progresivamente más amplia, hasta que toda la distribución de datos se aproxima a una distribución gaussiana con gran varianza. En el proceso inverso, la función score $\nabla_{\mathbf{x}} \log p_t(\mathbf{x})$ guía el retorno desde lo difuso hacia lo nítido.
\end{importantbox}

\subsection{Derivación de VE-SDE}

Se reescribe la cadena de Markov discreta en forma continua. Cuando el número de pasos de tiempo $N \to \infty$ y el intervalo temporal $\Delta t \to 0$, la probabilidad de transición se convierte en:

$$
\mathbf{x}_{t+\Delta t} = \mathbf{x}_t + \sqrt{\sigma^2(t+\Delta t) - \sigma^2(t)} \cdot \boldsymbol{\epsilon}, \quad \boldsymbol{\epsilon} \sim \mathcal{N}(0, \mathbf{I})
$$

Utilizando el desarrollo de Taylor $\sigma^2(t+\Delta t) \approx \sigma^2(t) + \frac{\mathrm{d}[\sigma^2(t)]}{\mathrm{d}t}\Delta t$ y tomando el límite, se obtiene:

$$
\mathrm{d}\mathbf{x} = \sqrt{\frac{\mathrm{d}[\sigma^2(t)]}{\mathrm{d}t}}\;\mathrm{d}\mathbf{w}
$$

Esta es la \textbf{SDE de varianza explosiva (VE-SDE)}. Los coeficientes correspondientes son:

\begin{itemize}
    \item \textbf{Coeficiente de arrastre}: $f(\mathbf{x}, t) = 0$ (no hay arrastre determinista)
    \item \textbf{Coeficiente de difusión}: $g(t) = \sqrt{\dfrac{\mathrm{d}[\sigma^2(t)]}{\mathrm{d}t}}$
\end{itemize}

\begin{knowledgebox}{¿Por qué se llama "varianza explosiva"}
En la VE-SDE, dado que el término de arrastre es cero, la media permanece constante ($\mathbb{E}[\mathbf{x}_t \mid \mathbf{x}_0] = \mathbf{x}_0$), mientras que la varianza $\sigma^2(t)$ aumenta monótonicamente con el tiempo y no tiene cota superior (si no se realiza un recorte). Por ello se denomina Variance Exploding: la varianza crece de manera "explosiva". La imagen física correspondiente es: cada punto de datos permanece "quieto en su sitio", pero la nube gaussiana de ruido que lo rodea se vuelve progresivamente más grande.
\end{knowledgebox}

\subsection{$\alpha_t$ y $\sigma_t$ en VE-SDE}

Bajo el marco unificado de Variational Diffusion Models, la VE-SDE corresponde a:

$$
\alpha_t = 1, \qquad \sigma_t = \sigma(t)
$$

Donde $\sigma(t)$ es la función de escala del ruido definida en SMLD. La relación señal-ruido (SNR) es:

$$
\mathrm{SNR}(t) = \frac{\alpha_t^2}{\sigma_t^2} = \frac{1}{\sigma^2(t)}
$$

Dado que $\sigma(t)$ aumenta monótonamente, la SNR disminuye monótonamente, cumpliendo las condiciones de Variational Diffusion Models.

\subsection{Resumen de la sección}

La VE-SDE es una extensión natural de SMLD en el dominio del tiempo continuo: sin arrastre, solo difusión, con una varianza que aumenta continuamente a lo largo del tiempo. Corresponde al proceso de aumento gradual de la varianza del ruido en las Annealed Langevin Dynamics.
\section{SDE de preservación de varianza: VP-SDE}

\subsection{De DDPM a VP-SDE}

La distribución de transición hacia adelante de DDPM se define como:

$$
q(\mathbf{x}_t \mid \mathbf{x}_{t-1}) = \mathcal{N}\!\left(\mathbf{x}_t;\, \sqrt{1-\beta_t}\,\mathbf{x}_{t-1},\, \beta_t \mathbf{I}\right)
$$

La distribución de salto correspondiente es:

$$
q(\mathbf{x}_t \mid \mathbf{x}_0) = \mathcal{N}\!\left(\mathbf{x}_t;\, \sqrt{\bar{\alpha}_t}\,\mathbf{x}_0,\, (1 - \bar{\alpha}_t)\mathbf{I}\right)
$$

Donde $\bar{\alpha}_t = \prod_{s=1}^{t}(1-\beta_s)$.

\begin{warningbox}{Diferencia clave entre DDPM y SMLD}
SMLD solo aumenta la varianza (el media permanece inalterado), mientras que DDPM simultáneamente \textbf{reduce el media} y aumenta la varianza —el factor $\sqrt{1-\beta_t}$ hace que el media decaiga gradualmente hacia cero. Esto significa que en DDPM los puntos de datos están al mismo tiempo "arrastrados hacia el origen" y "perturbados por ruido", convergiendo finalmente a una distribución gaussiana estándar $\mathcal{N}(0, \mathbf{I})$.
\end{warningbox}

\subsection{Derivación de VP-SDE}

Definiendo la varianza continua $\beta(t)$ (asociada mediante $\beta_t^{\text{discrete}} = \beta(t) \cdot \Delta t$), la cadena de Markov discreta de DDPM se convierte en:

$$
\mathbf{x}_{t+\Delta t} = \sqrt{1 - \beta(t)\Delta t}\;\mathbf{x}_t + \sqrt{\beta(t)\Delta t}\;\boldsymbol{\epsilon}
$$

Utilizando la aproximación de Taylor de primer orden $\sqrt{1 - \beta(t)\Delta t} \approx 1 - \frac{1}{2}\beta(t)\Delta t$ y tomando el límite $\Delta t \to 0$:

$$
\mathrm{d}\mathbf{x} = -\frac{1}{2}\beta(t)\,\mathbf{x}\,\mathrm{d}t + \sqrt{\beta(t)}\,\mathrm{d}\mathbf{w}
$$

Esta es la \textbf{VP-SDE} (Stochastic Differential Equation de Preservación de Varianza). Los coeficientes correspondientes son:

\begin{itemize}
    \item \textbf{Coeficiente de arrastre}: $f(\mathbf{x}, t) = -\frac{1}{2}\beta(t)\,\mathbf{x}$ (arrastre lineal que atrae a $\mathbf{x}$ hacia el origen)
    \item \textbf{Coeficiente de difusión}: $g(t) = \sqrt{\beta(t)}$
\end{itemize}

\begin{importantbox}{Origen del nombre VP-SDE}
Si la covarianza de los datos iniciales es la matriz identidad $\mathbf{I}$ (es decir, $\mathrm{Cov}(\mathbf{x}_0) = \mathbf{I}$), entonces bajo VP-SDE, para todos los momentos intermedios $t$, la covarianza permanece igual a $\mathbf{I}$:

$$
\mathrm{Cov}(\mathbf{x}_t) = \mathbf{I}, \quad \forall\, t
$$

Esto es el significado de "preservación de varianza" (Variance Preserving). El término de arrastre $-\frac{1}{2}\beta(t)\mathbf{x}$ en la ecuación difusiva compensa exactamente la varianza inyectada por el término de difusión, manteniendo la varianza total constante.
\end{importantbox}

\subsection{$\alpha_t$ y $\sigma_t$ en VP-SDE}

Bajo el marco de Variational Diffusion Models, VP-SDE corresponde a:

$$
\alpha_t = e^{-\frac{1}{2}\int_0^t \beta(s)\,\mathrm{d}s}, \qquad \sigma_t^2 = 1 - \alpha_t^2 = 1 - e^{-\int_0^t \beta(s)\,\mathrm{d}s}
$$

\begin{knowledgebox}{Relación entre $\bar{\alpha}_t$ discreto y $\alpha_t$ continuo}
En DDPM discreto, $\bar{\alpha}_t = \prod_{s=1}^t (1-\beta_s)$, lo cual es una forma de producto. En el límite continuo, se convierte en una forma exponencial $\alpha_t = e^{-\frac{1}{2}\int_0^t \beta(s)\,\mathrm{d}s}$. Ambos son aproximaciones equivalentes: utilizando $\ln(1-x) \approx -x$ (cuando $x$ es muy pequeño), $\prod(1-\beta_s) \approx e^{-\sum \beta_s} \approx e^{-\int \beta(s)\,\mathrm{d}s}$. Sin embargo, ten en cuenta que el exponente en el producto es $-\sum \beta_s$, mientras que el exponente en la versión continua contiene un factor $\frac{1}{2}$; esta diferencia proviene de los detalles de las aproximaciones $\sqrt{1-\beta}$ y $1-\frac{\beta}{2}$.
\end{knowledgebox}

\subsection{Monotonía decreciente del SNR}

La relación señal-ruido (SNR) en VP-SDE es:

$$
\mathrm{SNR}(t) = \frac{\alpha_t^2}{\sigma_t^2} = \frac{e^{-\int_0^t \beta(s)\,\mathrm{d}s}}{1 - e^{-\int_0^t \beta(s)\,\mathrm{d}s}}
$$

Dado que $\beta(t) > 0$, la integral $\int_0^t \beta(s)\,\mathrm{d}s$ crece estrictamente con respecto a $t$; por lo tanto, el numerador decrece monótonicamente y el denominador crece monótonicamente, resultando en que el SNR decrece monótonicamente. Esto valida que DDPM es un caso particular de Variational Diffusion Models.

\begin{warningbox}{$\beta(t)$ debe ser positivo}
En DDPM y VP-SDE, $\beta(t) > 0$ es una restricción fundamental. Si $\beta(t)$ fuera negativo en algún momento, el SNR podría dejar de decrecer monótonicamente y el proceso hacia adelante ya no sería un proceso de "adición de ruido" válido. En implementaciones prácticas, $\beta_t$ se define típicamente como un programa lineal o coseno que garantiza ser siempre positivo.
\end{warningbox}

\subsection{Resumen de la sección}

VP-SDE es una generalización de DDPM al dominio del tiempo continuo: el término de arrastre atrae los datos hacia el origen, mientras que el término de difusión inyecta ruido; ambos están cuidadosamente equilibrados para mantener la varianza constante. A diferencia de VE-SDE, VP-SDE posee un arrastre no nulo.
\section{SDE inversas y proceso generativo}

\subsection{Teorema de Anderson: SDE inversa}

Dado la SDE hacia adelante

$$
\mathrm{d}\mathbf{x} = f(\mathbf{x}, t)\,\mathrm{d}t + g(t)\,\mathrm{d}\mathbf{w}
$$

Anderson (1982) demostró que la SDE de tiempo inverso correspondiente es:

$$
\mathrm{d}\mathbf{x} = \left[f(\mathbf{x}, t) - g(t)^2 \nabla_{\mathbf{x}} \log p_t(\mathbf{x})\right]\mathrm{d}t + g(t)\,\mathrm{d}\bar{\mathbf{w}}
$$

donde $\bar{\mathbf{w}}$ es un proceso de Wiener en tiempo inverso, y $\mathrm{d}t$ representa aquí una \textbf{incremento negativo} de tiempo (el tiempo avanza desde $T$ hacia 0).

\begin{importantbox}{Estructura de la SDE inversa}
La SDE inversa comparte una estructura sorprendentemente similar a la SDE hacia adelante:
\begin{itemize}
    \item El coeficiente de difusión $g(t)$ es \textbf{idéntico};
    \item El término de arrastre resta \textbf{adicionalmente} $g(t)^2 \nabla_{\mathbf{x}} \log p_t(\mathbf{x})$ sobre la base del original $f(\mathbf{x}, t)$.
\end{itemize}
Lo único que se requiere conocer de manera adicional es la \textbf{función score} $\nabla_{\mathbf{x}} \log p_t(\mathbf{x})$, que es precisamente el objetivo aprendido por las redes neuronales en los modelos de difusión.
\end{importantbox}

\subsection{El papel de la función Score}

Repasando el enfoque utilizado en DDPM discreto:
\begin{enumerate}
    \item Definir el proceso hacia adelante (añadir ruido);
    \item Entrenar una red de predicción de ruido $\boldsymbol{\epsilon}_\theta(\mathbf{x}_t, t)$;
    \item Obtener la función score mediante la relación $\nabla_{\mathbf{x}} \log p_t(\mathbf{x}_t) = -\frac{\boldsymbol{\epsilon}}{\sigma_t}$;
    \item Ejecutar el proceso inverso de desruido utilizando la función score.
\end{enumerate}

En un marco de tiempo continuo, la lógica es totalmente consistente:
\begin{enumerate}
    \item Definir la SDE hacia adelante (arrastre + difusión);
    \item Entrenar una red score $\mathbf{s}_\theta(\mathbf{x}, t) \approx \nabla_{\mathbf{x}} \log p_t(\mathbf{x})$;
    \item Sustituir $\mathbf{s}_\theta$ en la SDE inversa para ejecutar la generación.
\end{enumerate}

\begin{knowledgebox}{Equivalencia entre score y predicción de ruido}
Para una distribución de salto $q(\mathbf{x}_t | \mathbf{x}_0) = \mathcal{N}(\alpha_t \mathbf{x}_0, \sigma_t^2 \mathbf{I})$, la función score puede expresarse como:
$$
\nabla_{\mathbf{x}_t} \log p_t(\mathbf{x}_t | \mathbf{x}_0) = -\frac{\mathbf{x}_t - \alpha_t \mathbf{x}_0}{\sigma_t^2} = -\frac{\boldsymbol{\epsilon}}{\sigma_t}
$$
donde $\boldsymbol{\epsilon}$ es el ruido gaussiano estándar muestreado al añadir ruido. Por lo tanto, predecir la función score es equivalente a predecir el ruido (la diferencia radica en dividir por $\sigma_t$). Esto explica por qué el objetivo de entrenamiento de "predicción de ruido" en DDPM y el objetivo de Score Matching son esencialmente lo mismo.
\end{knowledgebox}

\subsection{Implementación discreta del proceso inverso}

Aunque la SDE inversa es continua, se requiere discretizarla para el cálculo práctico. Dada una SDE inversa, dividimos el tiempo $[0, T]$ en $N$ pasos; en cada paso:

$$
\mathbf{x}_{t-\Delta t} \approx \mathbf{x}_t + \left[f(\mathbf{x}_t, t) - g(t)^2 \mathbf{s}_\theta(\mathbf{x}_t, t)\right](-\Delta t) + g(t)\sqrt{\Delta t}\;\boldsymbol{\epsilon}
$$

Esto es totalmente consistente con el proceso inverso discreto de DDPM: en cada paso se corrige la media utilizando la función score (o predicción de ruido) y luego se añade un ruido aleatorio apropiado.

\begin{warningbox}{Número de pasos de discretización y calidad de generación}
La discretización de una SDE es esencialmente el método de Euler-Maruyama. Cuantos más pasos, más precisa será la aproximación y mejor será la calidad de generación; sin embargo, también aumentará el costo computacional. DDPM requiere aproximadamente 1000 pasos para obtener una buena calidad, mientras que DDIM (la discretización correspondiente a una ODE determinista) solo necesita aproximadamente 50 pasos. En conferencias posteriores se presentarán solvers de SDE/ODE más eficientes (como DPM-Solver), lo cual permitirá reducir aún más el número de pasos a un rango de 10--15.
\end{warningbox}

\subsection{Resumen de la sección}

El teorema de Anderson garantiza que: dada cualquier SDE hacia adelante, la SDE inversa se determina automáticamente; lo único que requiere aprendizaje es la función score. Una vez discretizada la SDE inversa, recuperamos el algoritmo de muestreo de los modelos de difusión clásicos.

\section{Modelos de Difusión Variacional: un marco unificado}

\subsection{Forma general}

Los Modelos de Difusión Variacional definen una familia genérica de modelos de difusión. La forma general de la distribución de salto es:

$$
q(\mathbf{x}_t \mid \mathbf{x}_0) = \mathcal{N}\!\left(\mathbf{x}_t;\, \alpha_t \mathbf{x}_0,\, \sigma_t^2 \mathbf{I}\right)
$$

donde $\alpha_t$ y $\sigma_t$ son funciones diferenciables del tiempo $t$, que satisfacen:

\begin{enumerate}
    \item $\alpha_0 = 1$, $\sigma_0 = 0$ (en el momento inicial se trata de los datos originales);
    \item $\alpha_T \approx 0$, $\sigma_T \approx 1$ (en el momento final es aproximadamente una gaussiana estándar);
    \item \textbf{La relación señal-ruido SNR$(t) = \alpha_t^2/\sigma_t^2$ decrece monótonamente con respecto a $t$}.
\end{enumerate}

\begin{importantbox}{Significado de la disminución monótona del SNR}
El hecho de que el SNR decrezca monótonamente significa que, a medida que avanza el proceso hacia adelante, el componente de señal ($\alpha_t \mathbf{x}_0$) se debilita continuamente en relación con el componente de ruido ($\sigma_t \boldsymbol{\epsilon}$). Esta es la esencia del "añadido de ruido" en los modelos de difusión: la información se va sumergiendo progresivamente en el ruido, hasta que la distribución de los datos degenera en una distribución de puro ruido.
\end{importantbox}

\subsection{SMLD y DDPM como casos particulares}

\begin{table}[H]
\centering
\begin{tabular}{lccc}
\toprule
\textbf{Modelo} & $\alpha_t$ & $\sigma_t$ & \textbf{Tipo de SDE} \\
\midrule
SMLD / NCSN & $1$ & $\sigma(t)$ & VE-SDE \\
DDPM & $e^{-\frac{1}{2}\int_0^t \beta(s)\,\mathrm{d}s}$ & $\sqrt{1 - e^{-\int_0^t \beta(s)\,\mathrm{d}s}}$ & VP-SDE \\
\bottomrule
\end{tabular}
\caption{Parámetros de SMLD y DDPM bajo el marco de los Modelos de Difusión Variacional}
\end{table}

\begin{knowledgebox}{Apertura del espacio de diseño}
VE-SDE y VP-SDE son solo dos ejemplos concretos. Cualquier par $(\alpha_t, \sigma_t)$ que cumpla la condición de disminución monótona del SNR define un modelo de difusión válido. Esto ofrece a los investigadores un enorme espacio de diseño: pueden optimizar la calidad de generación, la velocidad de muestreo u otros indicadores eligiendo diferentes agendas de ruido. Por ejemplo, EDM (Elucidating the Design Space of Diffusion-Based Generative Models, Karras et al. 2022) exploró sistemáticamente este espacio de diseño.
\end{knowledgebox}

\subsection{De SDE a distribución de salto}

Dadas $f(\mathbf{x}, t)$ y $g(t)$ del SDE hacia adelante, se pueden deducir los correspondientes $\alpha_t$ y $\sigma_t$:

$$
\alpha_t = e^{\int_0^t f(s)\,\mathrm{d}s}, \qquad \sigma_t^2 = \alpha_t^2 \int_0^t \frac{g^2(s)}{\alpha_s^2}\,\mathrm{d}s
$$

A la inversa, dados $\alpha_t$ y $\sigma_t$, también se pueden recuperar $f(t)$ y $g(t)$:

$$
f(t) = \frac{\mathrm{d} \ln \alpha_t}{\mathrm{d}t}, \qquad g^2(t) = \frac{\mathrm{d}\sigma_t^2}{\mathrm{d}t} - 2\frac{\mathrm{d} \ln \alpha_t}{\mathrm{d}t}\sigma_t^2
$$

Esta conversión bidireccional implica que existen \textbf{tres descripciones equivalentes} —distribución de transferencia, distribución de salto y SDE— que pueden intercambiarse libremente.

\begin{warningbox}{Advertencia sobre confusión de símbolos}
Diferentes artículos utilizan convenciones de simbolismo distintas. Por ejemplo, el $\sigma$ de Song et al. se refiere a la agenda de ruido en VE-SDE (es decir, el $\sigma_i$ de SMLD), mientras que el $\sigma_t$ en los Modelos de Difusión Variacionales es la desviación estándar de la distribución de salto. En un mismo desarrollo pueden aparecer simultáneamente ambos tipos de $\sigma$, y deben distinguirse según el contexto. El profesor Minhyuk Sung también advirtió especialmente sobre este punto durante la clase.
\end{warningbox}

\subsection{Resumen de la sección}

Los Modelos de Difusión Variacional proporcionan un marco general mediante la parametrización $(\alpha_t, \sigma_t)$; SMLD y DDPM son sus casos particulares. Existe una correspondencia bidireccional exacta entre las tres descripciones (distribución de transferencia, distribución de salto y SDE).
\section{Análisis comparativo de VE-SDE y VP-SDE}

\subsection{Diferencias clave entre dos SDE}

Antes de profundizar en Probability Flow ODE, comparemos sistemáticamente las diferencias entre VE-SDE y VP-SDE; esto facilita la comprensión de los compromisos de diseño en distintos modelos de difusión.

\begin{table}[H]
\centering
\begin{tabular}{lcc}
\toprule
\textbf{Característica} & \textbf{VE-SDE (SMLD)} & \textbf{VP-SDE (DDPM)} \\
\midrule
Término de deriva $f(\mathbf{x},t)$ & $0$ & $-\frac{1}{2}\beta(t)\mathbf{x}$ \\
Coeficiente de difusión $g(t)$ & $\sqrt{\mathrm{d}[\sigma^2]/\mathrm{d}t}$ & $\sqrt{\beta(t)}$ \\
Atenuación de la media & No ($\alpha_t = 1$) & Sí ($\alpha_t < 1$) \\
Comportamiento de la varianza & Crecimiento monótono (exploding) & Mantenimiento constante (preserving) \\
Distribución de estado final & $\mathcal{N}(0, \sigma_T^2 \mathbf{I})$ & $\mathcal{N}(0, \mathbf{I})$ \\
Modo de atenuación del SNR & $1/\sigma^2(t)$ & $\alpha_t^2/(1-\alpha_t^2)$ \\
\bottomrule
\end{tabular}
\caption{Comparativa sistemática de VE-SDE y VP-SDE}
\end{table}

\begin{knowledgebox}{Intuición física}
La imagen física del \textbf{VE-SDE} es: cada punto de datos permanece "estacionario", pero la niebla de ruido que lo rodea se vuelve cada vez más densa; la varianza aumenta continuamente hasta ocultar por completo la estructura de los datos. La imagen física del \textbf{VP-SDE} corresponde a un proceso Ornstein-Uhlenbeck: los puntos de datos son atraídos simultáneamente hacia el origen por una fuerza elástica (atenuación de la media) y experimentan perturbaciones aleatorias (inyección de ruido); ambos factores se equilibran para mantener la varianza total constante.
\end{knowledgebox}

\subsection{Comparativa del rendimiento experimental}

Los experimentos de Song et al. (2021) muestran que, sobre el conjunto de datos CIFAR-10:
\begin{itemize}
    \item Al emplear muestreo DDPM con 1000 pasos, el FID (Frechet Inception Distance, donde menor es mejor) oscila alrededor de 3--4;
    \item Con DDIM, se alcanza un FID similar utilizando apenas unos 50 pasos;
    \item Mediante el uso de un solver SDE mejorado (por ejemplo, basado en métodos Predictor-Corrector), una calidad aceptable se logra con tan solo 10--15 pasos.
\end{itemize}

\begin{warningbox}{El FID no es la única métrica}
El FID mide simultáneamente la calidad generativa y la diversidad, pero no refleja completamente la percepción humana de la calidad de las imágenes. En aplicaciones prácticas, también deben considerarse la velocidad de generación, la cobertura de modos (diversity) y el desempeño en tareas posteriores. VE-SDE y VP-SDE presentan ventajas y desventajas relativas según la configuración utilizada.
\end{warningbox}

\subsection{Resumen de la sección}

VE-SDE y VP-SDE representan dos estrategias de ruido distintas: la primera implica difusión pura sin deriva, mientras que la segunda equilibra deriva y difusión. Comprender sus diferencias ayuda a tomar decisiones de diseño razonables en aplicaciones prácticas.
\section{EDO de Flujo de Probabilidad}

\subsection{De EDP a EDO}

La EDP estocástica hacia adelante describe un proceso estocástico cuyas trayectorias son en zigzag (ya que las trayectorias del proceso de Wiener, aunque continuas, no son diferenciables en ninguna parte). Una pregunta natural es:

\textit{¿Existe una EDO determinista tal que la distribución marginal $p_t(\mathbf{x})$ que genera sea exactamente la misma que la de la EDP estocástica?}

\begin{importantbox}{EDO de Flujo de Probabilidad}
Dada la EDP estocástica hacia adelante:
$$
\mathrm{d}\mathbf{x} = f(\mathbf{x}, t)\,\mathrm{d}t + g(t)\,\mathrm{d}\mathbf{w}
$$
Existe una única EDO (llamada \textbf{EDO de Flujo de Probabilidad}):
$$
\mathrm{d}\mathbf{x} = \left[f(\mathbf{x}, t) - \frac{1}{2}g(t)^2 \nabla_{\mathbf{x}} \log p_t(\mathbf{x})\right]\mathrm{d}t
$$
De modo que la distribución marginal en cada instante $t$ de las trayectorias de esta EDO coincide con la distribución marginal $p_t(\mathbf{x})$ de la EDP estocástica.
\end{importantbox}

Observe la relación entre la EDO de Flujo de Probabilidad y la EDP estocástica hacia atrás:

\begin{table}[H]
\centering
\begin{tabular}{lcc}
\toprule
& \textbf{Corrección del término de arrastre} & \textbf{Término estocástico} \\
\midrule
EDP estocástica hacia atrás & $f - g^2 \nabla_{\mathbf{x}} \log p_t$ & $g(t)\,\mathrm{d}\bar{\mathbf{w}}$ \\
EDO de Flujo de Probabilidad & $f - \frac{1}{2} g^2 \nabla_{\mathbf{x}} \log p_t$ & Ninguno \\
\bottomrule
\end{tabular}
\caption{Comparación entre la EDP estocástica hacia atrás y la EDO de Flujo de Probabilidad}
\end{table}

En la forma de la EDO, el coeficiente del score cambia de $g^2$ a $\frac{1}{2}g^2$, y al mismo tiempo se \textbf{elimina el término del proceso de Wiener}.

\subsection{Ventajas de las EDOs}

La EDO de Flujo de Probabilidad ofrece varias ventajas prácticas:

\begin{enumerate}
    \item \textbf{Muestreo determinista}: Dado el ruido inicial $\mathbf{x}_T$, la solución de la EDO es única; el mismo ruido siempre genera la misma imagen. Esto es consistente con el comportamiento de DDIM (de hecho, DDIM puede verse como una discretización de la EDO de Flujo de Probabilidad).
    \item \textbf{Menos pasos de muestreo}: Las trayectorias de las EDOs deterministas son suaves, lo que permite utilizar tamaños de paso más grandes y solucionadores numéricos de orden superior (como el método de Runge-Kutta), reduciendo significativamente el número de pasos de muestreo.
    \item \textbf{Cálculo preciso de la verosimilitud}: Mediante la fórmula del cambio instantáneo de variables, es posible calcular con precisión el logaritmo de la verosimilitud $\log p(\mathbf{x}_0)$ de los datos utilizando la EDO.
    \item \textbf{Interpolación semántica}: Dado que la EDO establece una biyección entre el espacio de datos y el espacio de ruido, es posible realizar interpolaciones en el espacio latente para obtener resultados intermedios con significado semántico.
\end{enumerate}

\begin{knowledgebox}{La relación entre DDIM y la EDO de Flujo de Probabilidad}
DDIM puede entenderse como un esquema de discretización especial de la EDO de Flujo de Probabilidad correspondiente a VP-EDP. Esto explica dos características centrales de DDIM: (1) el proceso de muestreo es determinista (no hay término estocástico), y (2) se puede generar con alta calidad utilizando muchos menos pasos que durante el entrenamiento (por ejemplo, 50 pasos en lugar de 1000).
\end{knowledgebox}

\subsection{Trayectorias de EDO vs. Trayectorias de EDP}

\begin{itemize}
    \item \textbf{Trayectoria de EDP}: En zigzag, estocástica; diferentes trayectorias que parten del mismo punto llegan a destinos distintos. La distribución marginal se obtiene promediando sobre todas las trayectorias posibles.
    \item \textbf{Trayectoria de EDO}: Suave, determinista; cada punto inicial corresponde a un único punto final. La distribución marginal está dada por la deformación global del conjunto de trayectorias (flujo de probabilidad).
\end{itemize}

Ambas distribuciones marginales $p_t(\mathbf{x})$ son exactamente iguales en cada instante $t$, pero las \textbf{trayectorias individuales} son diferentes.

\begin{warningbox}{La EDO de Flujo de Probabilidad sigue requiriendo el Score}
Aunque la EDO elimina el término estocástico, la función score $\nabla_{\mathbf{x}} \log p_t(\mathbf{x})$ aparece aún en el término de arrastre. Por lo tanto, el proceso de entrenamiento no cambia: ya sea que se utilice un muestreador basado en EDP o uno basado en EDO para generar muestras, el modelo sigue entrenando la función score (o equivalentemente, predicción del ruido).
\end{warningbox}

\subsection{Relación con Neural ODE y Normalizing Flows continuos}

La EDO de Flujo de Probabilidad revela las conexiones profundas entre los modelos de difusión, las Neural ODE (Chen et al., 2018) y los Continuous Normalizing Flows (CNF).

\begin{knowledgebox}{Continuous Normalizing Flows}
Los CNF definen una EDO $\mathrm{d}\mathbf{x}/\mathrm{d}t = v_\theta(\mathbf{x}, t)$ mediante un campo vectorial parametrizado por una red neuronal, mapeando distribuciones simples (como la gaussiana) a distribuciones complejas de datos. La EDO de Flujo de Probabilidad es exactamente de esta forma: su campo vectorial es $f(\mathbf{x}, t) - \frac{1}{2}g(t)^2 \nabla_{\mathbf{x}} \log p_t(\mathbf{x})$, donde la parte del score se aproxima mediante una red neuronal.

La diferencia clave radica en que los CNF entrenan directamente el campo vectorial, mientras que los modelos de difusión primero entrenan una red score mediante score matching y luego utilizan la EDO de Flujo de Probabilidad para el muestreo. El entrenamiento del segundo enfoque es más estable porque el score matching tiene un objetivo de desruido explícito.
\end{knowledgebox}

\begin{importantbox}{Cambio Instantáneo de Variables}
Para un flujo generado por una EDO $\mathrm{d}\mathbf{x}/\mathrm{d}t = \mathbf{v}(\mathbf{x}, t)$, la evolución del logaritmo de la densidad satisface:
$$
\frac{\mathrm{d}\log p_t(\mathbf{x}_t)}{\mathrm{d}t} = -\mathrm{tr}\left(\frac{\partial \mathbf{v}}{\partial \mathbf{x}}\right)
$$
Sustituyendo el campo vectorial de la EDO de Flujo de Probabilidad en esta fórmula, es posible calcular con precisión el logaritmo de la verosimilitud $\log p_0(\mathbf{x}_0)$ de los datos. Esto representa una de las ventajas importantes de los modelos de difusión frente a las GAN: estas últimas no pueden proporcionar valores exactos de verosimilitud.
\end{importantbox}

\subsection{Resumen de la sección}

La EDO de Flujo de Probabilidad es el "espejo" determinista de la EDP estocástica: conserva la misma distribución marginal, pero con trayectorias suaves y sin aleatoriedad. Esto abre las puertas al muestreo determinista, a soluciones eficientes y al cálculo preciso de verosimilitudes, estableciendo además conexiones profundas entre los modelos de difusión, las Neural ODE y los Continuous Normalizing Flows.
\section{Score Matching en tiempo continuo}

\subsection{Objetivo de entrenamiento}

En el marco de tiempo continuo, el objetivo de entrenamiento del score matching se puede expresar de manera unificada como:

$$
\mathcal{L} = \mathbb{E}_{t \sim \mathcal{U}(0,T)}\;\mathbb{E}_{\mathbf{x}_0 \sim p_{\text{data}}}\;\mathbb{E}_{\mathbf{x}_t \sim q(\mathbf{x}_t|\mathbf{x}_0)}\left[\lambda(t)\left\|\mathbf{s}_\theta(\mathbf{x}_t, t) - \nabla_{\mathbf{x}_t}\log q(\mathbf{x}_t|\mathbf{x}_0)\right\|^2\right]
$$

Donde:
\begin{itemize}
    \item $t$ se muestra de forma uniforme en $[0, T]$ (tiempo continuo);
    \item $\lambda(t)$ es una función de ponderación dependiente del tiempo;
    \item El score condicional $\nabla_{\mathbf{x}_t}\log q(\mathbf{x}_t|\mathbf{x}_0) = -\frac{\mathbf{x}_t - \alpha_t \mathbf{x}_0}{\sigma_t^2}$ tiene una forma analítica conocida.
\end{itemize}

\begin{importantbox}{Entrenamiento discreto vs. continuo}
En el DDPM discreto, el paso de tiempo $t$ se muestra uniformemente desde $\{1, 2, \ldots, T\}$. En el marco continuo, $t$ se muestra desde un intervalo continuo $[0, T]$. En la práctica del entrenamiento, ambas diferencias son prácticamente insignificantes: basta con reemplazar los índices de tiempo discretos por valores continuos de tiempo. La red $\mathbf{s}_\theta(\mathbf{x}, t)$ acepta $t$ continuo como entrada (normalmente procesado mediante codificación posicional sinusoidal).
\end{importantbox}

\subsection{Selección de la función de ponderación $\lambda(t)$}

Diferentes elecciones de $\lambda(t)$ corresponden a distintas estrategias de entrenamiento:

\begin{itemize}
    \item $\lambda(t) = g(t)^2$: corresponde al ponderado de verosimilitud en las EDE, relacionado con la optimización del límite variacional (VLB);
    \item $\lambda(t) = \sigma_t^2$: es la elección más común; equivale al objetivo de predicción de ruido $\|\boldsymbol{\epsilon}_\theta - \boldsymbol{\epsilon}\|^2$;
    \item $\lambda(t) = 1$: coincide directamente el score, pero puede ser inestable en momentos de bajo ruido (pequeños $t$).
\end{itemize}

\begin{knowledgebox}{De Score Matching a predicción de ruido}
Utilizando $\nabla_{\mathbf{x}_t}\log q(\mathbf{x}_t|\mathbf{x}_0) = -\boldsymbol{\epsilon}/\sigma_t$ y $\mathbf{s}_\theta = -\boldsymbol{\epsilon}_\theta / \sigma_t$, la pérdida de score matching se puede reescribir como:
$$
\left\|\mathbf{s}_\theta - \nabla_{\mathbf{x}_t}\log q\right\|^2 = \frac{1}{\sigma_t^2}\left\|\boldsymbol{\epsilon}_\theta - \boldsymbol{\epsilon}\right\|^2
$$
Por lo tanto, elegir $\lambda(t) = \sigma_t^2$ elimina exactamente el factor $1/\sigma_t^2$, obteniendo una pérdida de mínimos cuadrados medios (MSE) de predicción de ruido simplificada: este es precisamente el objetivo simplificado utilizado en el artículo original de DDPM.
\end{knowledgebox}

\subsection{Resumen de la sección}

El objetivo de entrenamiento del score matching en el marco de tiempo continuo es sustancialmente idéntico al del DDPM discreto; la única diferencia radica en que la variable de tiempo pasa de ser discreta a continua. La elección de la función de ponderación $\lambda(t)$ influye en la dinámica del entrenamiento, pero el objetivo más frecuente de predicción de ruido es completamente consistente desde ambas perspectivas.
\section{Síntesis y ampliación}

\subsection{Puntos clave de esta clase}

\begin{enumerate}
    \item \textbf{Marco unificado SDE}: Los modelos DDPM discretos y SMLD pueden generalizarse naturalmente a SDEs en tiempo continuo. El proceso hacia adelante queda completamente determinado por el coeficiente de deriva $f$ y el coeficiente de difusión $g$.

    \item \textbf{VE-SDE y VP-SDE}:
    \begin{itemize}
        \item VE-SDE (correspondiente a SMLD): $f = 0$, $g = \sqrt{\mathrm{d}[\sigma^2]/\mathrm{d}t}$, con varianza que aumenta monótonamente;
        \item VP-SDE (correspondiente a DDPM): $f = -\frac{1}{2}\beta(t)\mathbf{x}$, $g = \sqrt{\beta(t)}$, con varianza que se mantiene constante.
    \end{itemize}

    \item \textbf{SDE hacia atrás}: Se obtiene directamente del teorema de Anderson; solo es necesario conocer adicionalmente la función score. Una vez entrenada una red score, tanto el proceso hacia adelante como el hacia atrás quedan completamente determinados.

    \item \textbf{Modelos de Difusión Variacional}: A través de la parametrización $(\alpha_t, \sigma_t)$, bajo la condición de que la relación señal-ruido (SNR) disminuya monótonamente, se define una familia infinita de modelos de difusión. VE-SDE y VP-SDE son solo dos instancias de esta familia.

    \item \textbf{EDO de Flujo de Probabilidad}: Cada SDE tiene una EDO correspondiente que genera la misma distribución marginal pero con trayectorias deterministas. Esto proporciona herramientas para el muestreo eficiente y el cálculo de verosimilitud preciso.
\end{enumerate}

\begin{importantbox}{Valor de la perspectiva en tiempo continuo}
El valor del marco SDE en tiempo continuo no radica únicamente en "unificar" los métodos existentes. Lo más importante es que nos ofrece una metodología sistemática para diseñar nuevos modelos de difusión: basta con definir $f$ y $g$ (o equivalentemente, $\alpha_t$ y $\sigma_t$), mientras que lo restante (proceso hacia atrás, EDO, objetivo de entrenamiento) queda determinado automáticamente. Esta flexibilidad tipo "plug and play" ha impulsado una gran cantidad de trabajos posteriores.
\end{importantbox}

\subsection{Direcciones futuras}

El marco en tiempo continuo establecido en esta clase constituye la base teórica para varios temas importantes subsiguientes:

\begin{itemize}
    \item \textbf{Muestreadores eficientes}: Solucionadores numéricos de alto orden basados en la EDO de Flujo de Probabilidad (como DPM-Solver, DEIS, etc.) pueden generar muestras de alta calidad en solo 10--15 pasos;
    \item \textbf{Flow Matching / Rectified Flow}: Un método más conciso para entrenar modelos generativos en tiempo continuo, que puede considerarse una parametrización directa de la EDO de Flujo de Probabilidad;
    \item \textbf{Modelos de Consistencia}: Logran generación en un solo paso mediante la imposición de consistencia entre mapeos en diferentes momentos sobre la trayectoria de la EDO;
    \item \textbf{Evaluación de verosimilitud precisa}: Cálculo del logaritmo de la verosimilitud aprovechando la EDO y el estimador de traza de Hutchinson.
\end{itemize}

\subsection{Lecturas adicionales}

\begin{itemize}
    \item Song, Y., et al. (2021). ``Score-Based Generative Modeling through Stochastic Differential Equations.'' \textit{ICLR 2021.} —Referencia central de esta clase.
    \item Anderson, B. D. O. (1982). ``Reverse-time diffusion equation models.'' \textit{Stochastic Processes and their Applications.} —Teoría clásica de la SDE hacia atrás.
    \item Kingma, D. P., et al. (2021). ``Variational Diffusion Models.'' \textit{NeurIPS 2021.} —Propuesta del marco de Modelos de Difusión Variacional.
    \item Karras, T., et al. (2022). ``Elucidating the Design Space of Diffusion-Based Generative Models.'' \textit{NeurIPS 2022.} —Trabajo importante que explora sistemáticamente el espacio de diseño de modelos generativos basados en difusión en tiempo continuo.
\end{itemize}

%% --- Fin del contenido --- %%

\end{document}
